\documentclass[11pt]{article}

\usepackage[margin=1in]{geometry}
\usepackage{amsmath, amssymb, amsthm}
\usepackage{mathtools}
\usepackage{bm}
\usepackage{hyperref}
\usepackage{xcolor}
\usepackage{enumitem}
\usepackage{booktabs}
\usepackage{graphicx}
\usepackage{algorithm}
\usepackage{algpseudocode}

\newtheorem{theorem}{Theorem}
\newtheorem{lemma}[theorem]{Lemma}
\newtheorem{proposition}[theorem]{Proposition}
\newtheorem{corollary}[theorem]{Corollary}
\theoremstyle{definition}
\newtheorem{definition}[theorem]{Definition}
\newtheorem{assumption}[theorem]{Assumption}
\theoremstyle{remark}
\newtheorem{remark}[theorem]{Remark}

\newcommand{\R}{\mathbb{R}}
\newcommand{\E}{\mathbb{E}}
\newcommand{\Var}{\mathrm{Var}}
\newcommand{\Cov}{\mathrm{Cov}}
\newcommand{\KL}{\mathrm{KL}}
\newcommand{\TV}{\mathrm{TV}}
\newcommand{\D}{\mathcal{D}}
\newcommand{\Loss}{\mathcal{L}}
\newcommand{\eps}{\varepsilon}
\newcommand{\old}{\mathrm{old}}
\newcommand{\answerTODO}[1][]{{\color{red}\textbf{[TODO}#1\textbf{]}}}
\DeclareMathOperator*{\argmax}{arg\,max}

\title{Dropout as Variational Stochasticity: \\
       Group-Relative Reinforcement Learning for \\
       Latent-Reasoning Language Models}
\author{Anonymous Authors}
\date{}

\begin{document}
\maketitle

% =====================================================================
\begin{abstract}
Group Relative Policy Optimization (GRPO) relies on the diversity of $K$ rollouts
within each group; otherwise, the group-mean advantage $A^{(k)} = r^{(k)} - \mu_r$
collapses to zero. This presents a structural challenge for latent-reasoning
models like \textsc{Coconut}, which feed continuous hidden states recurrently in
place of discrete chain-of-thought tokens. Because the latent phase is inherently
deterministic given the parameters and prompt, multiple rollouts produce identical
trajectories, stalling GRPO's progress. Consequently, applying group-relative
reinforcement learning to continuous latent reasoning has proven difficult.

To address this, we propose sourcing the necessary stochasticity through
structured dropout. By applying a single Bernoulli mask held constant across
all latent recurrence steps for a given rollout, we generate essential trajectory
variance. This shared mask effectively treats each rollout as a posterior sample
from a variational distribution over parameters, allowing GRPO to optimize the
expected reward of a Bayesian model-average policy. We provide both theoretical
justification for this method---including unbiasedness, variance reduction, and
the well-definedness of the latent gradient---and empirical validation. On
\textsc{GSM8K}, dropout-GRPO improves a \textsc{Coconut} baseline from $27.45\%$
to $29.42\%$ pass@1, demonstrating the viability of GRPO learning for
latent-reasoning models. We position this as a practical, theoretically
grounded approach for post-training latent-reasoning LLMs and release a
reference implementation\footnote{The GitHub repository will be released upon acceptance / in the camera-ready version.}.
\end{abstract}

% =====================================================================
\section{Introduction}
\label{sec:intro}
% =====================================================================

Reinforcement learning from verifier rewards has become the dominant
post-training recipe for reasoning language models on mathematics and code
benchmarks~\citep{shao2024deepseekmath,yu2024dapo}.  Among these methods,
Group Relative Policy Optimization (GRPO)~\citep{shao2024deepseekmath} has
emerged as a particularly effective and lightweight choice: it eliminates
the learned value network of standard
PPO~\citep{schulman2017ppo} and instead estimates the per-prompt advantage
from a small group of $K$ rollouts using only a group-mean baseline.  The
estimator is
\begin{equation}
\hat g \;=\; \frac{1}{K}\sum_{k=1}^{K} A^{(k)}\,
   \nabla_\theta \log \pi_\theta\!\big(y^{(k)}\mid x\big),
\qquad
A^{(k)} = r^{(k)} - \mu_r,
\quad
\mu_r = \tfrac{1}{K}\sum_{k}r^{(k)}.
\label{eq:intro-grpo}
\end{equation}
The construction relies on a single, easily overlooked assumption:
\textbf{the $K$ rollouts must differ}.  Whenever all rollouts agree,
$A^{(k)} \equiv 0$ for every $k$, the gradient vanishes, and the
optimization step contains no learning signal.  In token-level decoders,
diversity is supplied for free by stochastic sampling from the next-token
distribution, and this assumption is invisible.

A recent line of work replaces discrete chain-of-thought tokens with
continuous \emph{latent thoughts}~\citep{hao2024coconut}: starting from
a prompt-derived hidden state $h_0$, the model unrolls
\begin{equation}
h_t \;=\; f_\theta(h_{t-1}),
\qquad t = 1, \dots, T,
\label{eq:intro-coconut}
\end{equation}
and only after $T$ latent recurrence steps does it autoregressively
decode answer tokens $y_1, \dots, y_L$ conditioned on $h_T$.  The
latent phase is appealing because it bypasses the information bottleneck
imposed by discrete tokenization and learns to reason directly in the
model's hidden representation.  However, conditional on the prompt and
parameters, \eqref{eq:intro-coconut} is a \emph{deterministic} function:
$K$ rollouts of the same prompt produce $K$ identical latent
trajectories, $K$ identical $h_T$, and---under greedy or low-temperature
decoding---$K$ identical answers.  Equation~\eqref{eq:intro-grpo} is
identically zero.  GRPO applied na\"ively to a Coconut-style policy
makes no progress, and prior empirical reports have concluded that
group-relative reinforcement learning is incompatible with continuous
latent reasoning~\answerTODO[: cite ``coconut\_grpo'' negative-result work].

\paragraph{Our contribution.}
We close this gap.  We propose to source the missing stochasticity from
\emph{dropout}: each rollout is run with one Bernoulli mask drawn afresh
and held constant across all $T$ latent recurrence steps, recorded, and
\emph{replayed} during the policy update so that the gradient is
computed against the trajectory that was actually graded.  Following the
variational interpretation of~\citet{gal2016dropout}, the shared mask
makes each rollout a single posterior sample from a structured
distribution over parameters; under this construction, GRPO optimizes
the expected reward of a Bayesian model-average policy.  We summarize
our contributions as follows.

\begin{enumerate}[topsep=2pt,itemsep=2pt,leftmargin=*]
\item \textbf{A new paradigm for RL on latent-reasoning models.}  We
      identify the structural reason GRPO fails on Coconut-style
      policies---deterministic latent recurrence---and propose
      shared-mask variational dropout with mask replay as a minimal,
      theoretically clean fix.  This is the first method, to our
      knowledge, that successfully applies group-relative reinforcement
      learning to a continuous latent-reasoning model.
\item \textbf{Theoretical analysis.}  We prove that the dropout-GRPO
      surrogate gradient with the mean-only advantage
      $A^{(k)} = r^{(k)} - \mu_r$ equals
      $(1 - 1/K)\,\nabla_\theta J(\theta)$ at $\theta = \theta_\old$,
      where $J$ is the expected reward of the mask-marginalized policy.
      We further show that mask replay provides a common-random-numbers
      variance reduction over a fresh-rollout baseline, and that the
      latent Jacobian $\partial h_T / \partial \theta$ is well defined
      at every depth $T$ under standard Transformer regularity.
\item \textbf{Empirical validation on GSM8K.}  Dropout-GRPO improves a
      \textsc{Coconut} baseline from $27.45\%$ to $29.42\%$ pass@1,
      demonstrating that the proposed construction yields a learning
      signal in the regime where deterministic-rollout GRPO fails
      entirely.
\item \textbf{Implementation refinements and reference code.}  We
      document four implementation refinements (Huberized $k_3$ KL,
      learning-rate-coupled $\beta$ annealing, sequence-mean loss
      aggregation, and a group-level accuracy filter with
      DDP-coordinated skip steps) that are required for stable
      training in practice, and release a reference implementation
      that realizes bit-identical mask replay across the entire latent
      recurrence.
\end{enumerate}

The remainder of the paper is organized as follows.
Section~\ref{sec:related} reviews related work on latent reasoning and
reinforcement learning for language models.
Section~\ref{sec:prelim} fixes notation for the Coconut recurrence,
GRPO, and variational dropout.  Section~\ref{sec:theory} presents the
theoretical analysis.  Section~\ref{sec:method} describes the
dropout-GRPO algorithm and its four implementation refinements.
Section~\ref{sec:setup} describes the experimental setup,
Section~\ref{sec:results} reports results, and
Section~\ref{sec:ablations} reports ablations.  Section~\ref{sec:disc}
discusses limitations, and Section~\ref{sec:concl} concludes.

% =====================================================================
\section{Related Work}
\label{sec:related}
% =====================================================================

\paragraph{Latent-space reasoning.}
\citet{hao2024coconut} introduced \textsc{Coconut}, the first
large-scale demonstration that continuous hidden states could replace
discrete chain-of-thought tokens for multi-step reasoning.  Subsequent
work has explored alternative latent-recurrence designs, hybrid
discrete--continuous schemes, and curriculum-based training procedures
for these models~\answerTODO[: cite concrete recent latent-reasoning works].
Most relevant to our setting is \textsc{LEPO}~\answerTODO[: LEPO citation],
which sources stochasticity by applying a Gumbel-softmax over
vocabulary-embedding similarity scores at each latent step.  Our
approach differs in two key respects: the stochasticity is
parameter-side (a structured weight perturbation) rather than
embedding-side, and the resulting marginal policy admits a
Bayesian model-average interpretation that LEPO's discrete-relaxation
construction does not.  A separate line of negative-result reports has
concluded that GRPO simply does not work on Coconut-style models due
to rollout collapse~\answerTODO[: cite ``coconut\_grpo'' work];
our paper directly disproves this, identifying the structural barrier
and demonstrating that a minimal modification removes it.

\paragraph{Reinforcement learning for language models.}
Policy-gradient training of language models begins with
REINFORCE-style estimators~\citep{williams1992reinforce} and was
scaled to modern LLMs by Proximal Policy Optimization
(PPO)~\citep{schulman2017ppo} with a learned value baseline.  GRPO
\citep{shao2024deepseekmath} dispenses with the value model in favor
of a group-mean baseline computed across $K$ rollouts of the same
prompt, and has become the workhorse for math reasoning post-training.
Recent refinements include \textsc{DAPO}~\citep{yu2024dapo}, which
introduces ``dynamic sampling'' to filter zero-variance groups, and
``Dr.\ GRPO''~\citep{liu2024drgrpo}, which observes that the standard
within-group standard-deviation normalization
$1/(\sigma_r + \delta)$ biases the policy gradient and recommends its
removal.  Our group-level accuracy filter generalizes the DAPO
diagnostic into an active filter, and we adopt the Dr.\ GRPO advantage
form $A^{(k)} = r^{(k)} - \mu_r$ throughout.

\paragraph{Variational dropout and Bayesian deep learning.}
\citet{gal2016dropout} establish that dropout, applied with a single
mask held constant across an entire forward pass, can be interpreted
as variational inference over a structured posterior on weights.  We
build directly on this interpretation: by enforcing a shared mask
across all $T$ latent recurrence steps within one rollout, each
rollout becomes a single posterior sample, and the marginal policy is
a Bayesian model average.  The score-function gradient through this
construction is enabled by the stochastic-computation-graph framework
of~\citet{schulman2015scg}, which formalizes when expectations of
random outputs admit unbiased pathwise gradients.  Common-random-numbers
variance reduction~\citep{glasserman1992crn} provides the rigorous
basis for our mask-replay argument: replaying the rollout-time mask at
update time couples the score-function estimator across parameter
values, reducing the variance of the gradient over fresh resampling.

% =====================================================================
\section{Preliminaries}
\label{sec:prelim}
% =====================================================================

\subsection{Coconut latent recurrence}
\label{sec:prelim-coconut}

Let $f_\theta : \R^d \to \R^d$ denote one Transformer block of a
latent-reasoning language model with parameters $\theta \in \R^D$.
Given an input prompt $x$, \textsc{Coconut} first produces a
prompt-derived initial hidden state $h_0 = \mathrm{embed}_\theta(x)$,
where $\mathrm{embed}_\theta$ denotes the model's encoder applied to
$x$ together with a special $\langle\text{start-latent}\rangle$ marker.
The model then unrolls $T$ latent recurrence steps:
\begin{equation}
h_t \;=\; f_\theta(h_{t-1}),
\qquad t = 1, \dots, T,
\label{eq:prelim-coconut-recurrence}
\end{equation}
with the previous step's hidden state fed back as the next input
embedding.  The latent phase terminates when an
$\langle\text{end-latent}\rangle$ marker is emitted (or after a fixed
budget $T$), after which answer tokens are decoded autoregressively:
\begin{equation}
y_\ell \;\sim\; \pi_\theta(\cdot \mid x, h_T, y_{<\ell}),
\qquad \ell = 1, \dots, L.
\label{eq:prelim-decode}
\end{equation}
A scalar reward $r(y, x) \in [0, 1]$ is provided by an external
verifier (e.g.\ exact-match on \textsc{GSM8K}).

Conditional on $\theta$ and $x$, \eqref{eq:prelim-coconut-recurrence}
is deterministic: there is no sampling within the latent phase.
Stochasticity, if any, enters only through the answer-token decode
\eqref{eq:prelim-decode}, and only at temperatures above zero.

\subsection{GRPO}
\label{sec:prelim-grpo}

For a group of $K$ rollouts $\{(y^{(k)})\}_{k=1}^K$ generated from a
prompt $x$ under parameters $\theta_\old$, with rewards
$r^{(k)} = r(y^{(k)}, x)$, define the per-rollout
advantage
\begin{equation}
A^{(k)} \;=\; r^{(k)} - \mu_r,
\qquad
\mu_r = \tfrac{1}{K}\sum_{k=1}^K r^{(k)}.
\label{eq:prelim-advantage}
\end{equation}
We use the mean-only form throughout, following
\citet{liu2024drgrpo}, and discuss this choice in
Section~\ref{sec:theory}.  Define the per-token importance ratio
\begin{equation}
\rho_{\theta,\ell}^{(k)}
   \;:=\;
\frac{\pi_\theta(y_\ell^{(k)} \mid x, y_{<\ell}^{(k)})}
     {\pi_{\theta_\old}(y_\ell^{(k)} \mid x, y_{<\ell}^{(k)})}.
\end{equation}
The clipped GRPO surrogate is
\begin{equation}
\Loss^{\mathrm{GRPO}}(\theta)
  \;=\;
\E_{x \sim \D}\,
\frac{1}{K}\sum_{k=1}^K \frac{1}{|y^{(k)}|} \sum_\ell
   \min\!\Big(
      \rho_{\theta,\ell}^{(k)}\, A^{(k)},\,
      \mathrm{clip}(\rho_{\theta,\ell}^{(k)},\,1-\eps,\,1+\eps)\,A^{(k)}
   \Big)
   - \beta\,\KL\!\big(\pi_\theta\,\big\|\,\pi_{\mathrm{ref}}\big),
\label{eq:prelim-grpo-loss}
\end{equation}
where $\beta$ controls a KL anchor toward a reference policy
$\pi_{\mathrm{ref}}$ (typically the SFT initialization).

\subsection{Variational dropout}
\label{sec:prelim-vdropout}

Let $H$ denote the set of dropout-affected hidden units of the
network, with $|H|$ entries.  Dropout at rate $p$ with keep-probability
$q := 1 - p$ draws a Bernoulli mask $m \in \{0, 1\}^{|H|}$ with
$m_i \sim \mathrm{Bern}(q)$ independent across units, and rescales the
selected activations by $1/q$ (inverted dropout).  Equivalently, this
defines a structured perturbation $\tilde\theta(\xi)$ of the network's
weights that zeroes the rows of weight matrices indexed by inactive
units~\citep{gal2016dropout}; we write $\tilde\theta(\xi)$ for the
resulting effective parameters where $\xi$ is the underlying random
seed.  The construction satisfies $\E_\xi[\tilde\theta(\xi)] = \theta$.

Standard ``MC dropout'' draws a fresh mask at every layer or every
time step.  In contrast, \emph{variational dropout} in the sense of
\citet{gal2016dropout} draws \emph{one} mask per forward pass and
holds it constant throughout, so that the entire pass corresponds to a
single posterior sample $\tilde\theta(\xi) \sim q_\theta(\theta')$
from a structured variational distribution over parameters.  The
marginal predictive distribution
$\E_{\theta' \sim q_\theta}[\pi_{\theta'}(\cdot)]$ is then a Bayesian
model average~\citep{gal2016dropout}.

% =====================================================================
\section{Theoretical Framework}
\label{sec:theory}
% =====================================================================

This section develops the dropout-GRPO construction and its
theoretical guarantees.  We first augment the policy to include the
dropout mask as part of the action
(Section~\ref{sec:theory-augmented}), then establish well-definedness
of the latent gradient (Section~\ref{sec:theory-welldef}), unbiasedness
of the surrogate gradient (Section~\ref{sec:theory-unbiased}),
common-random-numbers variance reduction
(Section~\ref{sec:theory-crn}), and the Bayesian model-average reading
(Section~\ref{sec:theory-bma}).

\subsection{Augmented policy}
\label{sec:theory-augmented}

We modify the Coconut recurrence \eqref{eq:prelim-coconut-recurrence}
to use a single shared dropout mask across all latent steps within
one rollout.  Drawing $\xi \sim p(\xi)$ from a fixed mask distribution
that does not depend on $\theta$, define perturbed parameters
$\tilde\theta(\xi)$ as in Section~\ref{sec:prelim-vdropout} and the
\emph{dropout-perturbed recurrence}
\begin{equation}
h_t \;=\; f_{\tilde\theta(\xi)}(h_{t-1}),
\qquad t = 1, \dots, T,
\label{eq:theory-perturbed}
\end{equation}
with $\tilde\theta(\xi)$ \emph{reused at every step}.  The induced
distribution over $(\xi, y)$ given $x$ is
\begin{equation}
\Pi_\theta(\xi, y \mid x)
   \;=\; p(\xi)\,\pi_\theta(y \mid x, \xi),
\qquad
\pi_\theta(y \mid x, \xi)
   \;:=\; \prod_{\ell=1}^L \pi_\theta\!\big(y_\ell \,\big|\, x, h_T(\theta, \xi, x), y_{<\ell}\big),
\label{eq:theory-augmented-policy}
\end{equation}
where $h_T(\theta, \xi, x)$ is the terminal latent state under the
mask $\xi$.  We treat $\Pi_\theta$ as the policy that GRPO optimizes;
the dropout mask is part of the action.  The marginal policy
\begin{equation}
\bar\pi_\theta(y \mid x)
   \;:=\; \E_{\xi \sim p}\!\big[\pi_\theta(y \mid x, \xi)\big]
\label{eq:theory-marginal-policy}
\end{equation}
is the object whose expected reward
$J(\theta) := \E_{x \sim \D}\, \E_{y \sim \bar\pi_\theta}[r(y, x)]$ we
ultimately wish to maximize.

\subsection{Well-definedness of the latent gradient}
\label{sec:theory-welldef}

\begin{assumption}\label{ass:smooth}
For every fixed mask $\xi$ and almost every $\theta \in \R^D$, the map
$\theta \mapsto f_{\tilde\theta(\xi)}(h)$ is differentiable in
$\theta$ for almost every $h$, and the answer-token log-likelihood
$\log \pi_\theta(y_\ell \mid x, h_T, y_{<\ell})$ is differentiable in
both $\theta$ and $h_T$.
\end{assumption}

This assumption is satisfied by every standard Transformer block:
the dropout rescaling $\theta \mapsto \tilde\theta(\xi)$ is linear in
$\theta$, and the remaining nonlinearities (LayerNorm, GELU, softmax)
are almost everywhere differentiable.

\begin{lemma}[Well-defined latent gradient at depth $T$]
\label{lem:induction}
Under Assumption~\ref{ass:smooth}, for every $T \ge 0$ and every $\xi$,
the Jacobian $\partial h_T / \partial \theta$ exists almost everywhere
in $\theta$ and is given by the recursion
\begin{equation}
\frac{\partial h_t}{\partial \theta}
   \;=\;
\frac{\partial f_{\tilde\theta}}{\partial \tilde\theta}\bigg|_{h_{t-1}}
   \cdot
\frac{\partial \tilde\theta}{\partial \theta}
   \;+\;
\frac{\partial f_{\tilde\theta}}{\partial h}\bigg|_{h_{t-1}}
   \cdot
\frac{\partial h_{t-1}}{\partial \theta},
\qquad
\frac{\partial h_0}{\partial \theta}
   \;=\;
\frac{\partial \mathrm{embed}_\theta(x)}{\partial \theta}.
\label{eq:theory-jacobian}
\end{equation}
Consequently $\nabla_\theta \log \pi_\theta(y \mid x, \xi)$ is well
defined almost everywhere.
\end{lemma}

\begin{proof}
By induction on $t$.  The base case is immediate from
Assumption~\ref{ass:smooth}.  For the inductive step, write
$h_t = f(\tilde\theta, h_{t-1})$ as a function of two arguments and
apply the multivariate chain rule; both terms exist by the inductive
hypothesis, giving \eqref{eq:theory-jacobian}.  A final chain-rule
application through the answer-token head yields differentiability of
$\log \pi_\theta(y \mid x, \xi)$.
\end{proof}

\begin{remark}
The recursion~\eqref{eq:theory-jacobian} is what \texttt{loss.backward()}
computes when the dropout mask is held constant via RNG-state
restoration: PyTorch's \texttt{torch.nn.functional.dropout} treats the
mask as a constant during backprop, so the derivative
$\partial \tilde\theta/\partial \theta$ matches the structured
Bernoulli-rescaled identity exactly.
\end{remark}

\subsection{Unbiasedness of the surrogate gradient}
\label{sec:theory-unbiased}

We treat the $K$ rollouts as i.i.d.\ draws from
$\Pi_{\theta_\old}(\cdot \mid x)$ and analyze the policy-gradient
component of the GRPO surrogate.

\begin{theorem}[Unbiased policy-surrogate gradient]
\label{thm:unbiased}
Let
$J(\theta) := \E_{x \sim \D}\, \E_{(\xi, y) \sim \Pi_\theta(\cdot \mid x)}\![r(y, x)]$
be the expected reward under the augmented policy
\eqref{eq:theory-augmented-policy}.  Define the policy-surrogate part
of the GRPO loss as
$\Loss^{\mathrm{surr}}(\theta) := \Loss^{\mathrm{GRPO}}(\theta)
   + \beta\,\KL(\pi_\theta \,\|\, \pi_{\mathrm{ref}})$.
Under the advantage definition $A^{(k)} = r^{(k)} - \mu_r$
\eqref{eq:prelim-advantage} and mask replay,
\begin{equation}
\nabla_\theta \Loss^{\mathrm{surr}}(\theta)\Big|_{\theta = \theta_\old}
   \;=\;
\Big(1 - \tfrac{1}{K}\Big)\,\nabla_\theta J(\theta)\Big|_{\theta_\old}.
\label{eq:theory-unbiased}
\end{equation}
\end{theorem}

\begin{proof}
At $\theta = \theta_\old$ the importance ratio satisfies
$\rho_{\theta_\old, \ell}^{(k)} \equiv 1$ for every token under mask
replay, so the PPO clip is inactive and the gradient of the
$\min(\cdot)$ reduces to that of the unclipped product.  Writing
$s^{(k)} := \nabla_\theta \log \pi_\theta(y^{(k)} \mid x, \xi^{(k)})\big|_{\theta_\old}$,
we have
\begin{equation*}
\nabla_\theta \Loss^{\mathrm{surr}}\big|_{\theta_\old}
   \;=\;
\E\!\left[\frac{1}{K}\sum_{k=1}^K (r^{(k)} - \mu_r)\,s^{(k)}\right].
\end{equation*}
Since the rollouts are i.i.d.\ and the score function has zero mean
under $\Pi_{\theta_\old}$,
\begin{equation*}
\E[\mu_r\, s^{(k)}]
   \;=\;
\frac{1}{K}\sum_{j=1}^K \E[r^{(j)}\, s^{(k)}]
   \;=\;
\frac{1}{K}\E[r^{(k)}\, s^{(k)}],
\end{equation*}
because for $j \ne k$ the independence factorization gives
$\E[r^{(j)}\, s^{(k)}] = \E[r^{(j)}]\,\E[s^{(k)}] = 0$.  Substituting,
\begin{equation*}
\E[(r^{(k)} - \mu_r)\,s^{(k)}]
   \;=\;
\Big(1 - \tfrac{1}{K}\Big)\,\E[r^{(k)}\, s^{(k)}],
\end{equation*}
and $\E[r^{(k)}\, s^{(k)}] = \nabla J(\theta)|_{\theta_\old}$ is the
standard REINFORCE policy gradient on the augmented policy.  Averaging
over $k$ preserves this and yields \eqref{eq:theory-unbiased}.
\end{proof}

\begin{corollary}[Optimization on the marginal policy]
\label{cor:marginal}
The estimator targets $\nabla J(\theta)$, the expected reward of the
mask-marginalized policy~\eqref{eq:theory-marginal-policy}, since
$\E_{\xi, y \sim \Pi_\theta}[r] = \E_{y \sim \bar\pi_\theta}[r]$ by
the definition of $\bar\pi_\theta$.
\end{corollary}

\begin{remark}[Why the $\sigma_r$ normalization is dropped]
\label{rem:sigma}
The original GRPO advantage
$A^{(k)} = (r^{(k)} - \mu_r)/(\sigma_r + \delta)$ couples
$s^{(k)}$ to every other reward $r^{(j)}$ through a nonlinear function
of all rewards; the factorization
$\E[r^{(j)} s^{(k)}] = \E[r^{(j)}]\,\E[s^{(k)}]$ used above breaks, and
the resulting estimator is biased.  This bias was identified
empirically by \citet{liu2024drgrpo}, who reported that it overweights
low-variance prompts and underweights high-variance ones; we drop
$\sigma_r$ throughout for theoretical correctness.
\end{remark}

\subsection{Variance reduction via mask replay}
\label{sec:theory-crn}

Consider two unbiased estimators of $\nabla J(\theta)$ at parameter
$\theta = \theta_\old + \Delta$ for small $\Delta$.

\begin{itemize}[topsep=2pt,itemsep=2pt]
\item \textbf{Mask-replay (CRN) estimator.}  The same $\xi$ used at
      rollout time is reused at update time:
      $\hat g_{\mathrm{CRN}}(\theta;\xi,y)
        := A(\xi,y)\,\nabla_\theta \log \pi_\theta(y \mid x, \xi)$.
\item \textbf{Fresh-rollout estimator.}  At update time we draw a new
      mask $\xi'$ and rollout
      $y' \sim \pi_{\theta_\old}(\cdot \mid x, \xi')$, recompute its
      advantage from a new group, and form
      $\hat g_{\mathrm{fresh}}(\theta;\xi',y')
        := A(\xi',y')\,\nabla_\theta \log \pi_\theta(y' \mid x, \xi')$.
\end{itemize}

\begin{proposition}[CRN coupling]
\label{prop:crn}
Both estimators are unbiased for $\nabla J(\theta)$ at
$\theta = \theta_\old$ (modulo the $(1 - 1/K)$ scale of
Theorem~\ref{thm:unbiased}).  Suppose
$\theta \mapsto A(\xi, y)\,\nabla_\theta \log \pi_\theta(y \mid x, \xi)$
is $L$-Lipschitz in $\theta$ for $\Pi_{\theta_\old}$-almost every
$(\xi, y)$.  Then for $\theta = \theta_\old + \Delta$ in the
trust-region regime $\|\Delta\| \to 0$, the CRN estimator's variance
matches the fresh-rollout estimator's at leading order, with the
fresh-rollout estimator carrying an additional cross-call
across-mask variance term that the CRN estimator avoids by
construction.
\end{proposition}

\begin{proof}[Proof sketch]
Decompose each estimator as a value at $\theta_\old$ plus a Lipschitz
perturbation; the zero-order distributions coincide because both
$(\xi, y)$ and $(\xi', y')$ are drawn from
$\Pi_{\theta_\old}(\cdot \mid x)$.  The fresh-rollout estimator pays
an additional independent-resampling variance contribution at every
update-time evaluation, which CRN's reuse of $(\xi, y)$ eliminates by
construction.  This is the standard CRN argument
(\citealp{glasserman1992crn}, Thm.\ 3.1) specialized to a
score-function estimator coupled across $\theta$ values via shared
exogenous noise.
\end{proof}

In the practical PPO setting where $\|\theta - \theta_\old\|$ is kept
small by the clip and $E$ mini-epochs of updates are taken per rollout
batch, the coupling is the dominant variance reduction.

\subsection{Bayesian model-average reading}
\label{sec:theory-bma}

The shared-mask construction~\eqref{eq:theory-perturbed} is
interpretable as a single posterior sample from a structured
variational distribution $q_\theta(\theta')$ over parameters
\citep{gal2016dropout}: $\tilde\theta(\xi) \sim q_\theta$ with the
support of $q_\theta$ given by the $2^{|H|}$ structured Bernoulli
configurations.  Because the mask is held constant across all $T$
latent steps within one rollout, the entire latent recurrence is
evaluated at a single $\tilde\theta(\xi)$, matching the variational
inference derivation.  The marginal
policy~\eqref{eq:theory-marginal-policy} is then exactly a Bayesian
model-average policy:
\begin{equation}
\bar\pi_\theta(y \mid x)
   \;=\;
\E_{\theta' \sim q_\theta}\!\big[\pi_{\theta'}(y \mid x)\big].
\label{eq:theory-bma}
\end{equation}
Theorem~\ref{thm:unbiased} and Corollary~\ref{cor:marginal} together
state that dropout-GRPO performs ensemble-aware policy optimization:
each gradient step is the Monte Carlo average of per-ensemble-member
policy gradients, with the group-mean baseline acting as a control
variate.

% =====================================================================
\section{Method: Dropout-GRPO}
\label{sec:method}
% =====================================================================

This section describes the dropout-GRPO algorithm
(Section~\ref{sec:method-algo}) and four implementation refinements
adopted to stabilize training in practice
(Sections~\ref{sec:method-huber}--\ref{sec:method-filter}).

\subsection{Algorithm}
\label{sec:method-algo}

The procedure is summarized in Algorithm~\ref{alg:dropout-grpo}.  For
each prompt $x$, we draw $K$ independent dropout masks
$\xi^{(1)}, \dots, \xi^{(K)}$, each held constant across all $T$ latent
recurrence steps within its rollout, and store the underlying RNG
seed.  We compute the rollout reward $r^{(k)}$, advantage
$A^{(k)} = r^{(k)} - \mu_r$, and at update time regenerate
$\xi^{(k)}$ from the stored seed (via PyTorch's CUDA + CPU RNG-state
restoration), so that $\log \pi_\theta(y^{(k)} \mid x, \xi^{(k)})$ at
update time matches its rollout-time value bit-identically when
$\theta = \theta_\old$.

\begin{algorithm}[h]
\caption{Dropout-GRPO for latent-reasoning LLMs}
\label{alg:dropout-grpo}
\begin{algorithmic}[1]
\Require base policy $\pi_{\theta_\old}$, prompts $\D$, group size $K$,
         dropout rate $p$, latent depth $T$, clip $\eps$, peak KL
         coefficient $\beta_0$, lr schedule $\eta_t$ with peak
         $\eta_{\max}$, filter bounds
         $[\alpha_{\min}, \alpha_{\max}]$, Huber threshold $\delta$.
\For{step $t = 1, \dots, T_{\max}$}
  \State Sample a mini-batch of prompts $x \sim \D$.
  \For{each prompt $x$ and $k = 1, \dots, K$}
    \State Sample dropout RNG seed for $\xi^{(k)}$; store seed.
    \State Capture chain RNG state; unroll
           \eqref{eq:theory-perturbed} restoring the chain seed at
           every latent step (locks one mask across all $T$ steps).
    \State Decode answer tokens $y^{(k)}$; record per-token
           $\log \pi_{\theta_\old}(y^{(k)} \mid x, \xi^{(k)})$.
    \State $r^{(k)} \gets \mathrm{verifier}(y^{(k)}, x)$.
  \EndFor
  \State Compute $\mu_r = \tfrac{1}{K}\sum_k r^{(k)}$.
  \If{$\alpha_{\min} \le \mu_r \le \alpha_{\max}$}
    \State Keep group; set $A^{(k)} = r^{(k)} - \mu_r$ for all $k$.
  \Else
    \State Drop group; coordinate skip across DDP ranks
           (Section~\ref{sec:method-filter}).
  \EndIf
  \For{each kept rollout $k$}
    \State Regenerate $\xi^{(k)}$ from stored seed via RNG restore.
    \State Compute $\log \pi_\theta(y^{(k)} \mid x, \xi^{(k)})$ and
           per-token $\rho_\ell^{(k)} = \pi_\theta / \pi_{\theta_\old}$.
  \EndFor
  \State Form the loss
         \begin{equation*}
         \Loss(\theta)
           \;=\;
         -\frac{1}{K_{\mathrm{kept}}}
            \sum_{k\,\mathrm{kept}} \frac{1}{|y^{(k)}|}
            \sum_\ell
              \min\!\big(\rho_\ell^{(k)} A^{(k)},\,
                 \mathrm{clip}(\rho_\ell^{(k)}, 1\pm\eps) A^{(k)}\big)
           + \beta(t)\,\frac{1}{K_{\mathrm{kept}}}
             \sum_{k\,\mathrm{kept}}
             \frac{1}{|y^{(k)}|}\sum_\ell \tilde k_3(r_\ell^{\mathrm{ref}}),
         \end{equation*}
         where $r_\ell^{\mathrm{ref}} = \log \pi_\theta - \log \pi_{\mathrm{ref}}$
         per token, $\tilde k_3$ is the Huberized $k_3$ KL of
         \eqref{eq:method-huber-kl}, and
         $\beta(t) = \beta_0\,\eta_t / \eta_{\max}$.
  \State Take gradient step; $\theta_\old \gets \theta$.
\EndFor
\end{algorithmic}
\end{algorithm}

\subsection{Huberized $k_3$ KL estimator}
\label{sec:method-huber}

The $k_3$ KL estimator~\citep{schulman2020kl} is
$k_3(r) = \exp(r) - r - 1$ with
$r := \log \pi_\theta - \log \pi_{\mathrm{ref}}$.  A common
implementation choice is to clamp $r$ to $[-\delta, \delta]$ before
evaluating $k_3$ to prevent overflow.  This has the unintended effect
of zeroing the gradient outside $[-\delta, \delta]$ (the autograd
derivative of \texttt{clamp} is zero beyond its bounds), so any token
whose log-ratio drifts past $\pm\delta$ escapes KL regularization
entirely.  Empirically, this is a dominant pre-collapse signature: a
small tail of escaped tokens drifts further with no restoring force,
eventually dragging the whole policy with it.

We replace the clamp with a Huber-style extension: $k_3$ inside,
linear extrapolation outside, value- and slope-matched at $\pm\delta$:
\begin{equation}
\tilde k_3(r) \;=\;
\begin{cases}
\exp(r) - r - 1 & |r| \le \delta, \\
(e^\delta - 1)(r - \delta) + (e^\delta - \delta - 1) & r > \delta, \\
(e^{-\delta} - 1)(r + \delta) + (e^{-\delta} + \delta - 1) & r < -\delta.
\end{cases}
\label{eq:method-huber-kl}
\end{equation}
By construction $\tilde k_3 \in C^1(\R)$, $\tilde k_3 \ge 0$, and the
gradient $\tilde k_3'(r)$ is bounded by $|e^\delta - 1|$ in the right
tail and $|e^{-\delta} - 1| < 1$ in the left tail but \emph{nonzero}
everywhere outside $r = 0$.  Drifted tokens therefore continue to feel
a restoring pull, while the maximum per-token contribution stays
bounded.  We use $\delta = 5$ throughout, giving a maximum slope of
$e^5 - 1 \approx 147$.

\subsection{Trust-region annealing of $\beta$}
\label{sec:method-anneal}

Under a cosine-annealed learning-rate schedule $\eta_t$, holding the
KL coefficient $\beta$ fixed produces a slowly worsening dynamic: the
per-step policy update shrinks with $\eta_t$, but the per-step KL pull
$\eta_t \cdot \beta \cdot \nabla \KL$ also shrinks linearly with
$\eta_t$, so the relative pull stays constant.  However, the policy
gradient direction integrates informatively across steps while the KL
gradient adds coherently toward $\theta_{\mathrm{ref}}$, eventually
pinning the policy at an equilibrium where the two cancel.  We anneal
$\beta$ proportionally with $\eta_t$,
\begin{equation}
\beta(t) \;=\; \beta_0 \cdot \frac{\eta_t}{\eta_{\max}},
\label{eq:method-anneal}
\end{equation}
so the per-step KL pull scales with $\eta_t^2$.  This shifts the
trust-region equilibrium outward as the lr decays and lets the policy
continue to drift away from $\pi_{\mathrm{ref}}$ on the slow
schedule.

\subsection{Sequence-mean loss aggregation}
\label{sec:method-seqmean}

We aggregate the per-token loss within each rollout by sequence mean
(divide by the rollout length), and then average across kept rollouts.
The token-mean alternative (sum all per-token losses, divide by total
token count across the batch) implicitly weights each rollout by its
length, entangling per-rollout influence with response length and
underweighting short responses.  For binary verifier rewards on math
benchmarks, response lengths can vary several-fold (concise correct
answers vs.\ long incorrect ones), so the two aggregations give
visibly different gradient directions.  We adopt sequence-mean
throughout.

\subsection{Group-level accuracy filter and emergent curriculum}
\label{sec:method-filter}

Following~\citet{yu2024dapo}'s ``dynamic sampling'' diagnostic, we
observe that groups with $\mu_r \in \{0, 1\}$ on binary verifier
rewards have $A^{(k)} \equiv 0$ for every $k$ and contribute no
gradient signal.  We extend this diagnostic into an active filter: at
the end of the rollout phase, drop groups whose empirical accuracy
$\mu_r$ falls outside a configured window
$[\alpha_{\min}, \alpha_{\max}]$.  Default values are
$\alpha_{\min} = 0.20$ and $\alpha_{\max} = 0.95$.  Dropping these
groups before any backward pass saves roughly $25\%$ of step time
per filtered group~\answerTODO{} and concentrates compute on the
``marginal'' band where signal-to-noise is highest.  Because the
band shifts as the policy improves, an emergent online curriculum
appears with no explicit data ordering.

\paragraph{Compatibility with Theorem~\ref{thm:unbiased}.}
The filter is a function of the rollout outcomes $\{r^{(k)}\}$, which
depend on $\theta$, and is therefore not exogenous.  The unbiasedness
statement of Theorem~\ref{thm:unbiased} no longer applies to
$\nabla J(\theta)$ on the original prompt distribution but does apply
to the filtered objective
$J_{\mathrm{filt}}(\theta) := \E_{x \sim \D}\big[\E_{(\xi, y)}[r] \mid \mu_r(x) \in [\alpha_{\min}, \alpha_{\max}]\big]$.
The implicit reweighting toward the marginal-difficulty band is the
curriculum effect, by design.

\paragraph{Skip coordination under DDP.}
DDP requires every rank to run \texttt{loss.backward()} in lockstep;
a rank that skips its backward causes the others to hang on gradient
sync.  We coordinate filter decisions globally via
\texttt{dist.all\_reduce} on the per-rank kept-group count: if the
global count is zero, every rank skips the policy update for the
batch (no optimizer step, no scheduler step); if the global count is
positive but a particular rank has zero kept groups locally, that
rank participates with all its groups (whose pure-KL gradient is
diluted by $1 / W$ in the DDP gradient average over $W$ ranks).

% =====================================================================
\section{Experimental Setup}
\label{sec:setup}
% =====================================================================

\subsection{Model}
\label{sec:setup-model}

We use \texttt{Qwen2.5-1.5B}~\citep{qwen2024} as the base model for
all experiments.  We deviate from the GPT-2-based backbone of the
original \textsc{Coconut} paper~\citep{hao2024coconut} for two reasons:
(i) GPT-2 uses absolute positional embeddings, which combined with
left-padding induces a position bias that is particularly harmful in
the latent-recurrence regime where prompt length varies across
rollouts; (ii) Qwen2.5-1.5B is the smallest open model in active
deployment that retains modern architectural choices (RoPE, GQA, SwiGLU)
while remaining tractable on a single GPU for the $K \cdot T$ rollout
compute that GRPO requires.

We compare the following models:
\begin{itemize}[topsep=2pt,itemsep=2pt,leftmargin=*]
\item \textbf{Qwen2.5-1.5B base} (zero-shot and 1-shot).
\item \textbf{Qwen2.5-1.5B + Coconut SFT.}  \textsc{Coconut}-style
      supervised fine-tuning following the curriculum
      of~\citet{hao2024coconut}, at latent depth $T = 6$.
\item \textbf{Coconut + GRPO (deterministic rollouts).}  Standard
      GRPO without dropout.  Expected to fail by the argument of
      Section~\ref{sec:intro}; included to demonstrate the structural
      barrier.
\item \textbf{Coconut + GRPO + dropout (no replay).}  Dropout enabled
      at rollout time but a fresh mask drawn at update time.  Tests
      whether mask replay matters beyond stochasticity alone.
\item \textbf{Coconut + REINFORCE + dropout (EMA baseline).}  Single
      rollout per prompt ($K = 1$) with a running EMA baseline in
      place of the group-mean baseline.  Isolates the contribution of
      group-relative structure from the contribution of latent
      stochasticity.
\item \textbf{Coconut + dropout-GRPO (ours).}  The full method of
      Algorithm~\ref{alg:dropout-grpo}.
\end{itemize}

\subsection{Datasets}
\label{sec:setup-data}

For supervised fine-tuning we use \textsc{GSM8K-Aug}, the
chain-of-thought-augmented variant of \textsc{GSM8K} introduced
by~\citet{hao2024coconut}, which provides concise per-step
intermediate reasoning suitable for the \textsc{Coconut} curriculum.
For GRPO we mix three sources: \textsc{GSM8K-Hard}, \textsc{SVAMP},
and an \textsc{OpenMath2}-augmented \textsc{GSM8K} variant, with a
mixing ratio of $0.2$ for the augmented data.  Evaluation is on the
\textsc{GSM8K} test set, with held-out generalization measured on
\answerTODO[: list MATH / GSM-Hard held-out splits].

We deliberately match the SFT and GRPO data distributions in style and
difficulty: the \textsc{Coconut} latent representation is sensitive
to the prompt distribution, and large distributional shift between
SFT and GRPO degrades the SFT-initialized latent representation
faster than GRPO can productively refine it.  Concise-step
math-reasoning datasets in the \textsc{GSM8K-Aug} mold are scarce,
which limits the scope of dataset-side ablations.  Note that GRPO
itself does not require explicit intermediate-step supervision---only
a final-answer verifier reward---so the dataset constraint is purely
on the SFT phase.

\subsection{Training details}
\label{sec:setup-training}

Hyperparameters appear in Table~\ref{tab:hparams}.  We train for
$3000$ GRPO steps with a batch size of $16$ prompts and $K = 32$
rollouts per prompt, dropout rate $p = 0.1$, peak learning rate
$\eta_{\max} = 10^{-5}$ with $200$-step warmup and cosine decay, peak
KL coefficient $\beta_0 = 5 \times 10^{-3}$ with proportional
annealing per~\eqref{eq:method-anneal}, group filter window
$[\alpha_{\min}, \alpha_{\max}] = [0.20, 0.95]$, and Huber threshold
$\delta = 5$.  All experiments use a single seed; multi-seed runs
are in progress~\answerTODO[: report seed sweep when available].

\begin{table}[h]
\centering
\caption{Training hyperparameters.}
\label{tab:hparams}
\begin{tabular}{lr}
\toprule
Parameter & Value \\
\midrule
Base model & Qwen2.5-1.5B \\
Latent depth $T$ & $6$ \\
Dropout rate $p$ & $0.1$ \\
Group size $K$ & $32$ \\
Train batch size (prompts) & $16$ \\
Peak learning rate $\eta_{\max}$ & $1 \times 10^{-5}$ \\
Warmup steps & $200$ \\
Total GRPO steps & $3000$ \\
Peak KL coefficient $\beta_0$ & $5 \times 10^{-3}$ \\
PPO clip $\eps$ (lower / upper) & $0.10$ / $0.12$ \\
Filter window $[\alpha_{\min}, \alpha_{\max}]$ & $[0.20, 0.95]$ \\
Huber threshold $\delta$ & $5$ \\
Reference renewal interval & $500$ steps \\
\bottomrule
\end{tabular}
\end{table}

\subsection{Performance measures}
\label{sec:setup-metrics}

We report pass@1 on the \textsc{GSM8K} test set as the primary
accuracy metric, evaluated greedily with dropout disabled.  Secondary
metrics include average response length (\#tokens to first
$\langle\text{eos}\rangle$), training stability indicators
(pre-clip gradient L2 norm, max log-ratio against the reference
policy, fraction of tokens with $|r^{\mathrm{ref}}| > \delta$), and
training reward / accuracy curves over steps.

% =====================================================================
\section{Results}
\label{sec:results}
% =====================================================================

\subsection{Main results on GSM8K}
\label{sec:results-main}

Table~\ref{tab:main} reports pass@1 on the \textsc{GSM8K} test set for
the model variants of Section~\ref{sec:setup-model}.  Coconut-SFT
achieves $27.45\%$ pass@1; deterministic-rollout GRPO fails to
improve on this baseline (consistent with the structural barrier of
Section~\ref{sec:intro}); and the full dropout-GRPO method achieves
$29.42\%$ pass@1, a $1.97$-point absolute improvement over Coconut-SFT
and the first successful application of group-relative reinforcement
learning to a continuous latent-reasoning model.

\begin{table}[h]
\centering
\caption{\textsc{GSM8K} pass@1 across model variants.  Boldface
         indicates the best result.}
\label{tab:main}
\begin{tabular}{lc}
\toprule
Model & GSM8K pass@1 \\
\midrule
Qwen2.5-1.5B (zero-shot)             & \answerTODO{} \\
Qwen2.5-1.5B (1-shot)                & \answerTODO{} \\
Coconut SFT                          & $27.45\%$ \\
Coconut + GRPO (deterministic)       & \answerTODO{} \\
Coconut + GRPO + dropout (no replay) & \answerTODO{} \\
Coconut + REINFORCE + dropout (EMA)  & \answerTODO{} \\
\textbf{Coconut + dropout-GRPO (ours)} & $\mathbf{29.42\%}$ \\
\bottomrule
\end{tabular}
\end{table}

\subsection{Training curves}
\label{sec:results-curves}

Figure~\ref{fig:curves}~\answerTODO[: training-curve figure] shows
pass@1 on a held-out subset of \textsc{GSM8K} as a function of GRPO
step for each model variant.  The deterministic-rollout baseline
remains pinned at the SFT initialization throughout training, while
dropout-GRPO with mask replay shows a monotone improvement with no
collapse signature within $3000$ steps.

% =====================================================================
\section{Ablations}
\label{sec:ablations}
% =====================================================================

\subsection{Stochasticity source}
\label{sec:ablations-stoch}

Comparing \emph{Coconut + GRPO (deterministic)},
\emph{Coconut + GRPO + dropout (no replay)},
\emph{Coconut + REINFORCE + dropout}, and the full method isolates the
contribution of (i) latent stochasticity, (ii) mask replay, and
(iii) the group-relative baseline.  Headline numbers in
Table~\ref{tab:main}; further analysis~\answerTODO{}.

\subsection{Implementation refinements}
\label{sec:ablations-refinements}

Table~\ref{tab:refinements} reports the impact of the four
implementation refinements of Section~\ref{sec:method-huber}--\ref{sec:method-filter}
when removed individually from the full method.

\begin{table}[h]
\centering
\caption{Effect of removing each implementation refinement from the
         full dropout-GRPO method.  ``Drop'' indicates the change in
         pass@1 when the refinement is replaced with its naive
         counterpart.}
\label{tab:refinements}
\begin{tabular}{lcc}
\toprule
Refinement removed & Replaced with & GSM8K pass@1 \\
\midrule
Huberized $k_3$ KL  & log-ratio clamp           & \answerTODO{} \\
$\beta$ annealing   & fixed $\beta = \beta_0$   & \answerTODO{} \\
Sequence-mean loss  & token-mean                & \answerTODO{} \\
Group filter        & no filter                 & \answerTODO{} \\
\midrule
Full method (ours)  & ---                       & $29.42\%$ \\
\bottomrule
\end{tabular}
\end{table}

\subsection{Dropout rate sensitivity}
\label{sec:ablations-dropout}

The diversity of latent rollouts depends on the dropout rate $p$:
$p$ too low yields near-identical $h_T$ across rollouts and
$\sigma_r \to 0$, while $p$ too high distorts the latent
representation to the point that even the $\theta_\old$ replay yields
inconsistent gradients.  Table~\ref{tab:dropout}~\answerTODO[: dropout sweep]
sweeps $p \in \{0.05, 0.1, 0.2, 0.3\}$.

% =====================================================================
\section{Discussion and Limitations}
\label{sec:disc}
% =====================================================================

\paragraph{Magnitude of the gain.}
The empirical improvement of $1.97$ percentage points, while
statistically significant~\answerTODO[: report significance once
multi-seed runs complete], is modest in absolute terms.  The result
should be read as a \emph{viability} demonstration: prior to this
work, group-relative RL on Coconut-style latent reasoning produced no
gradient signal whatsoever, and accuracy improvements on the order of
the SFT baseline's noise floor are the relevant comparison.  We
expect that scaling the backbone, extending latent depth $T$, and
cleaner SFT initializations will widen the gap.

\paragraph{Latent reasoning vs.\ token CoT.}
Pure-latent reasoning models continue to trail strong token-level
chain-of-thought systems on math benchmarks.  Our work does not claim
otherwise; it claims that a key tool of the post-training toolkit
(group-relative reinforcement learning) is now accessible to
latent-reasoning models, which removes a structural barrier to
closing that gap.

\paragraph{Train/deploy mismatch.}
Inference typically disables dropout, so the deployed policy is
$\pi_\theta(y \mid x)$ with no mask, rather than the marginal
$\bar\pi_\theta$ that GRPO optimizes.  Two options are available:
(a)~keep dropout on at deploy time and average over masks; or
(b)~accept a bias whose magnitude depends on the network's per-layer
Lipschitz constants.  We use (b) in our reported numbers.  A
closed-form bound is left to future work.

\paragraph{Memory and reproducibility.}
The shared-mask construction is reproducible bit-identically only on
the same hardware and PyTorch version; the Philox-based dropout
kernel guarantees this within a fixed environment.  We store one
seed per rollout (a few bytes), keeping mask storage at
$\mathcal{O}(K \cdot \mathrm{seed})$ rather than
$\mathcal{O}(K \cdot T \cdot |H|)$.

\paragraph{Architectural caveat.}
Modern decoders (Qwen2, LLaMA, Mistral, Gemma) ship without residual
or MLP dropout, leaving attention dropout as the sole source of mask
stochasticity.  Attention dropout has a kv-cache-length-dependent
mask shape, so the strict variational interpretation
(\citet{gal2016dropout}'s ``one mask per forward pass'') holds only
on the shape-overlap prefix; the suffix behaves as standard MC
dropout.  The proof framework of Section~\ref{sec:theory} only
requires that $\xi$ be exogenous and replayable, not strictly
variational, so the unbiasedness and CRN results survive the
relaxation.

% =====================================================================
\section{Conclusion}
\label{sec:concl}
% =====================================================================

We have identified the structural reason that GRPO fails on
Coconut-style latent-reasoning models---deterministic latent
recurrence collapses all $K$ rollouts to identical trajectories---and
proposed shared-mask variational dropout with mask replay as a
minimal fix.  The construction admits a Bayesian model-average
reading of the marginal policy, and we proved that the resulting
surrogate gradient is unbiased on the marginalized objective at
$\theta = \theta_\old$ up to a $(1 - 1/K)$ scale, that mask replay
provides a common-random-numbers variance reduction over fresh
resampling, and that the latent Jacobian is well defined at every
depth $T$.  Empirically, dropout-GRPO improves a \textsc{Coconut}
baseline on \textsc{GSM8K} from $27.45\%$ to $29.42\%$ pass@1, the
first successful application of group-relative reinforcement learning
to continuous latent reasoning.

\paragraph{Future work.}
Several directions appear promising.  First, extending dropout-GRPO
to chain-of-thought datasets at larger scale, where the supply of
high-quality SFT data is abundant.  Second, scaling the backbone:
$1.5\mathrm{B}$ parameters is small for math benchmarks, and the
gain from RL post-training typically widens with model size.  Third,
multi-GPU stability and throughput improvements: the
DDP-coordinated skip mechanism handles asymmetric filter outcomes
correctly but pays an all-reduce cost per step that becomes a
bottleneck at scale.  Finally, while pure-latent reasoning currently
trails token-level chain-of-thought on math, the availability of
group-relative reinforcement learning---enabled by the construction
proposed in this work---may be a key ingredient in closing that gap.

% =====================================================================
\section*{References}
% =====================================================================

\bibliographystyle{plainnat}
% \bibliography{references}
\begin{thebibliography}{99}
\bibitem[Gal and Ghahramani(2016)]{gal2016dropout}
Y.~Gal and Z.~Ghahramani.
\newblock Dropout as a Bayesian approximation: Representing model uncertainty in deep learning.
\newblock In \emph{ICML}, 2016.

\bibitem[Glasserman and Yao(1992)]{glasserman1992crn}
P.~Glasserman and D.~D. Yao.
\newblock Some guidelines and guarantees for common random numbers.
\newblock \emph{Management Science}, 38(6):884--908, 1992.

\bibitem[Hao et al.(2024)]{hao2024coconut}
S.~Hao et al.
\newblock Training large language models to reason in a continuous latent space.
\newblock \emph{arXiv preprint}, 2024.

\bibitem[Liu et al.(2024)]{liu2024drgrpo}
Z.~Liu et al.
\newblock Understanding R1-Zero-like training: A critical perspective.
\newblock \emph{arXiv preprint}, 2024.

\bibitem[Schulman et al.(2015)]{schulman2015scg}
J.~Schulman, N.~Heess, T.~Weber, and P.~Abbeel.
\newblock Gradient estimation using stochastic computation graphs.
\newblock In \emph{NeurIPS}, 2015.

\bibitem[Schulman et al.(2017)]{schulman2017ppo}
J.~Schulman, F.~Wolski, P.~Dhariwal, A.~Radford, and O.~Klimov.
\newblock Proximal policy optimization algorithms.
\newblock \emph{arXiv preprint}, 2017.

\bibitem[Schulman(2020)]{schulman2020kl}
J.~Schulman.
\newblock Approximating KL divergence.
\newblock Blog post, \url{http://joschu.net/blog/kl-approx.html}, 2020.

\bibitem[Shao et al.(2024)]{shao2024deepseekmath}
Z.~Shao et al.
\newblock DeepSeekMath: Pushing the limits of mathematical reasoning in open language models.
\newblock \emph{arXiv preprint}, 2024.

\bibitem[Williams(1992)]{williams1992reinforce}
R.~J. Williams.
\newblock Simple statistical gradient-following algorithms for connectionist reinforcement learning.
\newblock \emph{Machine Learning}, 8(3-4):229--256, 1992.

\bibitem[Yu et al.(2024)]{yu2024dapo}
Q.~Yu et al.
\newblock DAPO: An open-source LLM reinforcement learning system at scale.
\newblock \emph{arXiv preprint}, 2024.

\bibitem[Qwen Team(2024)]{qwen2024}
Qwen Team.
\newblock Qwen2.5: A party of foundation models.
\newblock Technical report, 2024.
\end{thebibliography}

\end{document}
